% Testing and Evaluation Section

\section{Testing and Evaluation}

\subsection{Our Testing Philosophy}
We didn't wait until the end to start testing. Instead, we built checks into the system at every level, from individual functions up to the full end-to-end flow.

\textbf{Unit tests} catch problems at the smallest scale. Every data cleaning function, every feature extraction routine, and every model wrapper has tests that verify it does what it claims to do. We used pytest for this since it's straightforward and widely supported.

\textbf{Integration tests} make sure the pieces work together. The preprocessing pipeline feeds into model training, which feeds into API endpoints—and we tested those handoffs to make sure data doesn't get lost or mangled along the way.

\textbf{End-to-end smoke tests} walk through the entire flow, from raw data ingestion all the way to charts rendering in the browser. These aren't exhaustive, but they catch the obvious breakages before anyone else sees them.

\textbf{Load tests} gave us a rough sense of performance. Nothing fancy—just enough to confirm that the API responds in reasonable time and doesn't fall over under moderate traffic.

\textbf{Regression tests} run automatically whenever code changes. If someone accidentally breaks a model or an API contract, the tests catch it before the change gets merged.

\subsection{Specific Things We Tested}
A few examples of what our test suite covers:
\begin{itemize}
\item Feeding the data ingestion module malformed records—missing fields, bad timestamps, weird characters—to make sure it handles garbage gracefully instead of crashing.
\item Running the NLP feature extraction on edge cases like super-short reviews, reviews with emoji, and the occasional non-English snippet that slips through.
\item Training a model, saving it to disk, loading it back, and verifying it still makes the same predictions. Sounds obvious, but serialization bugs can be sneaky.
\item Checking that the API returns exactly the JSON structure the frontend expects, including proper error messages when something goes wrong.
\item Rendering dashboard charts with no data, with huge datasets, and with multiple overlapping series to make sure the UI holds up under different conditions.
\end{itemize}

\subsection{How We Measured Success}
We tracked several types of metrics:

For the \textbf{forecasting models themselves}, we measured Mean Absolute Error (how far off predictions typically land), Root Mean Squared Error (which penalizes big misses more heavily), and Mean Absolute Percentage Error (to compare accuracy across features with different scales). We used holdout validation and rolling-origin testing to get honest estimates of how well models would perform on future data.

For the \textbf{system as a whole}, we monitored API response times (specifically the 95th percentile), success rates, and memory usage during model inference.

For \textbf{code quality}, we tracked test coverage on critical modules and counted how many bugs the CI pipeline caught before they reached production.

\subsection{Automation and CI}
Every pull request triggers the full test suite through GitHub Actions. Unit tests, integration tests, and even a quick model validation step all run automatically. If the model validation fails—say, a code change unexpectedly tanks prediction accuracy—the PR can't merge without explicit reviewer approval. This keeps the main branch stable and trustworthy.

\subsection{Testing Results}
In the end, we achieved solid coverage across the core pipeline. Unit tests cover the data processing and model components, integration tests verify the API contracts, and smoke tests confirm the frontend renders correctly. We prioritized making the system reproducible and reliable over chasing every last edge case. For a research prototype, that's the right tradeoff.
